\SourceLeaf{216}%
\bigskip\Needspace{10\baselineskip}
\pdfbookmark[1]{Further Report. Ceres Ferdinandea. May 1802}{ceres180205}
\begin{center}
{\small FURTHER REPORT ON THE NEW MAJOR PLANET.\par}
CERES FERDINANDEA.\par
\medskip\hrule\medskip
Monatliche Correspondenz. Vol. V. pp.\ 462—481. May 1802.\par
\medskip\hrule
\end{center}\medskip
—\quad —\quad —\quad Through 3 March our observations agreed with \textsc{Gauss}'s elements, as we reported in the previous issue, p.\ 389 [p.\ 209 of this volume]. Since then Dr.\ \textsc{Gauss} has continued the comparison, as shown below.
\begin{center}
{\small\setlength{\tabcolsep}{5pt}
\begin{tabular}{r r|r|r|r|r}
\multicolumn{2}{c|}{\shortstack{Seeberg\\1802}}&\shortstack{Calculated\\R.A. of $\CeresSymbol$}&Difference&\shortstack{Calculated\\decl. of $\CeresSymbol$}&Difference\\\hline
March&6&$185^0\ 49'\ 23''3$&$+4''6$&$16^0\quad4'\ 12''2$&$+23''0$\\
—&7&$185\ 38\ 22.0$&$+6.1$&$16\ 10\ 45.6$&$+29.7$\\
—&10&$185\quad3\ 57.2$&$+8.3$&$16\ 29\ 48.4$&$+29.5$\\
—&11&$184\ 52\quad5.2$&$+9.4$&$16\ 35\ 55.2$&$\ .\ .\ .\ .\ .$\\
—&15&$184\quad3\quad2.9$&$+10.9$&$16\ 58\ 58.2$&$+27.5$\\
—&16&$183\ 50\ 28.4$&$+9.4$&$17\quad4\ 20.6$&$+21.7$\\
—&17&$183\ 37\ 48.5$&$+13.7$&$17\quad9\ 32.6$&$+26.2$\\
—&18&$183\ 25\quad3.2$&$+12.5$&$17\ 14\ 34.4$&$+29.1$\\
—&19&$183\ 12\ 14.4$&$+15.8$&$17\ 19\ 25.1$&$+34.3$
\end{tabular}}
\end{center}

Dr.\ \textsc{Gauss} also compared two \textit{Greenwich} observations with his elements, and the right ascensions agree extremely closely with ours if one accepts the assumption stated above, which is beyond doubt, that the minute in the observation of 19 February was written down incorrectly by oversight.
\begin{center}
{\small\setlength{\tabcolsep}{2.4pt}
\begin{tabular}{l r|r|r|r|r|r|r}
\multicolumn{2}{c|}{}&&\multicolumn{2}{c|}{Difference for}&&\multicolumn{2}{c}{Difference for}\\
\multicolumn{2}{c|}{\shortstack{Greenwich\\1802}}&\shortstack{Calculated\\R.A. of $\CeresSymbol$}&\shortstack{Dr.\ \textsc{Mas-}\\\textsc{kelyne}}&\textsc{v. Zach}&\shortstack{Calculated\\declination\\of $\CeresSymbol$}&\shortstack{Dr.\ \textsc{Mas-}\\\textsc{kelyne}}&\textsc{v. Zach}\\\hline
February&19&$187^0\ 58'\ 13''3$&$-4''3$&$-4''1$&$14^0\ 20'\ 36''2$&$+35''3$&$+21''1$\\
March&6&$185\ 49\quad3.8$&$+5.0$&$+4.6$&$16\quad4\ 24.0$&$+34.5$&$+23.0$
\end{tabular}}
\end{center}

This comparison now clearly indicates that observations of the planet will depart increasingly from \textsc{Gauss}'s ellipse toward the end of March. Yet even if they were off by 5 minutes in June, which is unlikely, that would be of \textit{no consequence at all}. The error can only build up gradually, so by always comparing the preceding observations with the elements or with the ephemeris, one will always know it in advance with. sufficient precision. Dr.\ \textsc{Gauss} also writes to us that he has already made preparations to adjust the elements again at once so that they match the latest observations as closely as possible, should the error build up to one minute. One can then expect very close agreement again for a considerable time. But as soon as the observations for this year are finished and gathered together, he intends to determine the true ellipse with all its perturbation equations as precisely as possible.
\begin{flushright}[\textsc{von Zach}.]\end{flushright}
